\documentclass[10pt,twocolumn]{article}
\usepackage[margin=0.72in]{geometry}
\usepackage{amsmath,amssymb,amsthm}
\usepackage{booktabs}
\usepackage{microtype}
\usepackage{algorithm}
\usepackage[noend]{algpseudocode}
\usepackage[hidelinks]{hyperref}
\usepackage{xcolor}

\newcommand{\cone}{\operatorname{cone}}
\newcommand{\conv}{\operatorname{conv}}
\newcommand{\erank}{\operatorname{erank}}
\newcommand{\R}{\mathbb{R}}
\newcommand{\bx}{\mathbf{x}}
\newcommand{\bq}{\mathbf{q}}
\newcommand{\bv}{\mathbf{v}}
\newcommand{\bw}{\mathbf{w}}
\newcommand{\bmu}{\boldsymbol{\mu}}
\newcommand{\bV}{\mathbf{V}}
\newcommand{\bW}{\mathbf{W}}
\newcommand{\bX}{\mathbf{X}}
\newcommand{\bS}{\mathbf{S}}
\newcommand{\good}[1]{\textcolor{black}{\textbf{#1}}}

\title{\bf The Best Possible Raw Conic Classifier\\for Class-Incremental Learning}
\author{Khiem Dinh Phi}
\date{}

\begin{document}
\maketitle

\begin{abstract}
A conic read-out models each class as a polyhedral cone $C_c=\{\bV_c^{\!\top}\bw:\bw\ge0\}$
and scores a query by $\lVert\Pi_{C_c}\bq\rVert$. It is an attractive
class-incremental primitive: it is fitted per class, so it never revisits a past
image; it is a strict generalisation of multi-prototype matching; and unlike a
discriminative head it produces an absolute, not relative, membership score.
This report asks how far a \emph{raw} conic classifier---no fusion with a
discriminative head, no post-hoc calibration---can be pushed, and settles the
attribution for every design axis in the family. On ImageNet-R $T{=}10$ the best
raw cone reaches $\mathbf{80.07}$ A-Last against a RanPAC bar of $80.28$, a
deficit of $0.21$, from $71.02$ for a plain nearest-class-mean. The gain
decomposes as: the tied within-class metric $+6.65$, the conic rule over
whitened prototypes $+0.25$, \emph{discriminative} generator selection (oriented
PCA against foreign material) $+1.58$, and size-adaptive per-class ray budgets
$+0.57$. The value of non-negativity is
\emph{atom-dependent} and its mechanism is tail-driven: at k-means atoms an
unconstrained rank-$R$ subspace ties the cone ($78.02$ vs $77.92$), but at the
discriminative atoms the method actually uses the cone wins by $+0.60$
($79.77$ vs $79.17$). Averaged over all classes the constraint is
\emph{anti}-discriminative---it shaves $0.0245$ off the true class against
$0.0144$ off a random foreign one---yet against the single \emph{top competing}
class, the pair that decides, the sign reverses to $+0.0074$ ($+6.6\%$ of the
decision margin), rescuing $196$ top-1 contests and breaking $91$. Two further
findings are negative. (i) The residual deficit is a capacity mismatch, not
geometry: binned by class size, the raw cone beats RanPAC by $+14.23$ on classes
with under 50 training rows and loses by $-6.24$ on classes with over 200.
(ii) The cone already wins outright in the small-sample regime---$+1.86$ A-Last
on ImageNet-A at $R{=}4$.
We give the full recipe, the ledger of what each component is worth, and the
list of mechanisms that were tried and failed.
\end{abstract}

\section{Setting and protocol}

A ViT-B/16 encoder $\phi:\mathcal{X}\to\R^{768}$
(\texttt{vit\_base\_patch16\_224.augreg\_in21k}) is adapted with LoRA on the
first task only and then frozen; all later tasks contribute features and
statistics but no gradients. This is the A\textsuperscript{+} first-session
adaptation recipe. Write $\bx=\phi(\cdot)/\lVert\phi(\cdot)\rVert$. For each
class $c$ the read-out stores a small set of vectors $\bV_c$ and produces a
score $s_c(\bq)$; prediction is $\arg\max_{c\in\text{seen}}s_c(\bq)$.

Because the backbone is frozen after task $0$, per-class objects fitted at a
class's birth stage remain valid for the rest of the run, and final-stage
accuracy over all classes \emph{is} A-Last exactly. Read-out ablations are
therefore roughly $10\times$ cheaper than a full staged replay, which is what
makes the sweep in this report affordable.

Unless stated otherwise every number is ImageNet-R, $T{=}10$, seed $0$, A-Last
in per cent, with RanPAC on the same features as the bar ($80.28$ A-Last,
$85.13$ A-Avg). \textbf{Seeds are discussed in \S\ref{sec:limits} and are the
main threat to everything below.}

\paragraph{What ``raw'' excludes.} No fusion with RanPAC, no per-class score
calibration, no test-time adaptation. This restriction is deliberate: $\beta$
selection in a cone/RanPAC fusion was measured to swing a headline by $0.77$
under a $0.09$ perturbation of the features, which is larger than most effects
studied here. Fusion results are reported in \S\ref{sec:fusion} as a ceiling,
not as the method.

\section{The object and its score}

\subsection{One rule, four names}
Given unit generators $\bV=[\bv_1,\dots,\bv_R]^{\!\top}$, the conic hull is
\begin{equation}
C(\bV)=\Big\{\textstyle\sum_r\lambda_r\bv_r:\lambda_r\ge0\Big\}
      =\{\bV^{\!\top}\bw:\bw\ge0\}.
\end{equation}
For unit $\bq$, Moreau's decomposition gives
$\langle\bq,\Pi_C\bq\rangle=\lVert\Pi_C\bq\rVert^2$, hence
\begin{equation}
s(\bq)=\cos\angle(\bq,\Pi_C\bq)=\lVert\Pi_C\bq\rVert,
\label{eq:score}
\end{equation}
the exact analogue of a subspace read-out's $\lVert P\bq\rVert$ with the
projection taken onto a cone. Equation~\eqref{eq:score} implies that the
``residual coverage'' score $\cos\cdot\lVert\mathrm{rec}\rVert$ equals
$\cos^2$, and that the ``angular'' score $1-\tfrac{2}{\pi}\arccos\cos$ and the
``NNLS residual'' score $1-\tfrac12\sqrt{2-2\cos}$ are monotone in $\cos$. All
four induce identical rankings. Empirically they agree to four decimal places
(\texttt{exp17}: $88.36/88.36/88.36$ on CIFAR-100 $T{=}10$; only the
covariance-reweighted variant differs, at $88.18$). Treating them as separate
arms measures one quantity repeatedly.

\subsection{Nesting: the cone contains multi-prototype}
Restricting $\bw$ to its $k$ largest-cosine coordinates gives
$C_1\subseteq C_2\subseteq\cdots\subseteq C_R$ with $s_k$ nondecreasing in $k$.
At $k{=}1$ the score reduces \emph{exactly} to $\max_r\langle\bq,\bv_r\rangle$.
The conic hull is therefore a one-parameter generalisation of multi-prototype
matching with hard nearest-prototype as its degenerate member. This identity is
also the cheapest available correctness assertion, and it matters: our NNLS
solver once ran one of five hundred FISTA iterations for its entire history
because a convergence test compared a tensor with itself. Every conic score
produced before the fix was the one-step approximation
$\bw=\eta\,\mathrm{relu}(\bV\bq)$. On one cell the corrected solver moves
near-OOD AUROC from $0.622$ to $0.847$ (\texttt{exp20}, iteration sweep
$1/2/5/25/100/500 \to 0.622/0.676/0.789/0.846/0.847/0.847$), and it reverses the
sign of the ray-budget trend. \emph{At $k_{\text{local}}{=}1$ the conic score
must equal max-cosine bit-for-bit}; assert it.

\subsection{Storage and what is legal}
The read-out stores $R_c$ unit vectors per class plus one shared $d\times d$
within-class scatter. At $R{=}32$ that is $32\cdot768\cdot200=4.9$M floats
($19.7$\,MB), against RanPAC's $10000\cdot200=2.0$M ($8$\,MB): the cone already
has $2.5\times$ more parameters and still loses. \textbf{The deficit is not
capacity in the parameter-count sense; it is what those parameters are fitted
to.} No past image is ever revisited. When a class is born it may see: its own
rows, the rows of the other classes arriving in the same task, and the stored
generators of every past class.

\section{The design space}

Six axes are separable and were measured separately. In rough order of what
they turned out to be worth:
\begin{enumerate}\itemsep1pt
\item \textbf{Metric} --- the space generators are fitted and scored in.
\item \textbf{Generator selection} --- which $R$ directions represent a class.
\item \textbf{Discrimination} --- whether selection sees other classes.
\item \textbf{Allocation} --- how many rays each class gets.
\item \textbf{Constraint} --- $\bw\ge0$ versus alternatives.
\item \textbf{Calibration} --- correcting cone-size bias post hoc.
\end{enumerate}

\subsection{Axis 1: the metric ($+6.65$)}
Let $\Sigma$ be the pooled within-class scatter. It is \emph{additive across
classes} and therefore accumulates task by task with no stored samples, exactly
as RanPAC's Gram matrix does. With $\bW=\operatorname{chol}(\Sigma^{-1})$ we fit
and score in $\bx\mapsto\bW^{\!\top}\bx/\lVert\cdot\rVert$. Generators are built
in the birth-time whitened space, mapped back through $\bW_{\text{birth}}^{-1}$
and stored there; at any later stage they are re-metricated as
$\widehat{\bV\bW_t}$. Storage stays at $R_c$ vectors per class, no past sample
is revisited, and the approximation is empirically lossless---an oracle that
refits past classes from raw rows gains $+0.05$.

This axis alone lifts nearest-class-mean from $71.02$ to $77.67$. It is by a
wide margin the largest lever, and it is not conic at all.

Its regularisation is saturated. Sweeping the shrinkage $\delta$ in
$\Sigma+\delta\,\tfrac{\operatorname{tr}\Sigma}{d}I$ over four orders of
magnitude moves the cone by $0.45$ ($79.32$--$79.77$); Ledoit--Wolf
independently selects $\delta{=}0.0094$ and $\delta{=}0.1$ scores within $0.03$
of it (\texttt{exp34}: $79.67$ vs $79.70$). There are no points left in this
knob. The LDA extension---adding the between-class scatter
$\bS_b=\sum_c n_c(\bmu_c-\bar\bmu)(\bmu_c-\bar\bmu)^{\!\top}$, which is also
computable from stored means with zero images---is structurally starved in CIL:
$\operatorname{rank}\bS_b\le n_{\text{seen}}-1$, so at stage $0$ with 20 classes
it offers at most 19 dimensions.

\subsection{Axis 2: generator selection ($+0.25$, but $+61.8$ against the wrong choice)}

Both classical choices are the same constrained factorisation with different
solutions. Write $\bV=\mathbf{G}\bX$ with $\mathbf{G}\ge0$,
$\mathbf{G}\mathbf{1}=\mathbf{1}$:
\begin{itemize}\itemsep1pt
\item \textbf{SPA / separable NMF.} $\mathbf{g}_r=\mathbf{e}_{i_r}$: each
generator \emph{is} a data point, chosen greedily by maximum residual norm
(equivalently QR with column pivoting). Under separability this provably
recovers the true extremal rays.
\item \textbf{$k$-means / convex NMF.} $\mathbf{g}_r$ is the uniform
distribution on a cluster: each generator is a convex combination of data
points.
\end{itemize}
SPA sits at a vertex of the constraint set and $k$-means at a barycentre; they
optimise outward extent versus central representativeness. A centroid averages
$n/R$ samples and has variance $\sigma^2R/n$; an SPA ray is a single
sample---and specifically the argmax over $n$ residual norms, so its noise lies
in the upper tail. At $n\approx100$, $R{=}4$ the variance ratio exceeds
$25\times$ before the selection bias.

\begin{table}[t]\centering\small
\begin{tabular}{lcc}
\toprule
generators & classifier & near-OOD \\
\midrule
random ($R$ data points) & 54.85 & 0.806 \\
SPA (extremal)           & 36.15 & 0.756 \\
$k$-means (centroids)    & \good{74.25} & 0.844 \\
convex-NMF (learned)     & 73.80 & 0.841 \\
archetypal               & 73.97 & 0.840 \\
\bottomrule
\end{tabular}
\caption{Generator choice, ImageNet-R, $R{=}4$, matched budget
(\texttt{exp26}). SPA loses to \emph{random} by 18.7 points. Convex-NMF drives
reconstruction error 7.5\% \emph{below} its $k$-means initialisation and still
loses the task: reconstruction is the wrong objective.}
\label{tab:gens}
\end{table}

Table~\ref{tab:gens} is the largest measured effect anywhere in this line of
work ($36.15\to74.25$ at $R{=}4$; $11.60\to73.43$ at $R{=}2$). The lesson is not
that $k$-means is good---it is that \emph{containment is the wrong target}.
Adding a generator can only enlarge the cone, so a containment-optimal set
maximises coverage, which is precisely the property that destroys discrimination
when the score is a distance to that set. SPA optimises coverage; $k$-means does
not. Three learned dictionaries (convex-NMF, archetypal analysis, and a
gradient-fitted variant) all land within $0.5$ of $k$-means, which says the
reconstruction objective is exhausted.

\paragraph{Naming.} The resulting object is still a conic hull---$C(\bV)$ is
closed under positive scaling, convex, polyhedral and pointed---but it is the
conic hull \emph{of the dictionary}, not of the class:
$C(\bV_{k\text{-means}})\subsetneq\conv\text{-}\cone(\bX)$ and
$\bX\not\subseteq C(\bV_{k\text{-means}})$. Non-negativity of the coefficients
is preserved; containment of the class is not, and neither is the separability
recovery guarantee. ``A polyhedral cone spanned by $R$ learned generators'' is
the honest phrase; ``conic hull'' unqualified will be read as $\cone(\bX)$ and
imports containment intuitions that no longer hold.

\subsection{Axis 3: discrimination in the construction ($+1.58$)}
\label{sec:opca}

Both incumbents are \emph{blind to other classes}. A per-class generative fit
cannot make its score class-selective, which is exactly the observed deficit:
$k$-means cone accuracy falls from $79.50$ at $R{=}4$ to $77.92$ at $R{=}32$,
below a plain whitened class mean, because adding rays raises $s_c$ for foreign
queries as fast as for own-class ones. The fix is to put foreign information
into the construction. The objective is
\begin{equation}
\max_{\bV}\ \sum_{\bx\in X_c}\lVert\Pi_{C(\bV)}\bx\rVert^2
\;-\;\gamma\sum_{\mathbf{a}\in F}\lVert\Pi_{C(\bV)}\mathbf{a}\rVert^2 ,
\label{eq:disc}
\end{equation}
where $F$ is the foreign material legally available at class $c$'s birth: the
fit rows of the other classes in the current task plus the stored rays of every
past class, capped at $F_{\max}{=}2000$ by random subsample. Never a past image.

Since $\mathbb{E}\lVert P\bq\rVert^2=\operatorname{tr}(P^{\!\top}\bS P)$,
relaxing $\Pi_C$ to an orthogonal projector makes \eqref{eq:disc} exactly the
top-$R$ eigenproblem of $\bS_c-\gamma\bS_F$. This is \textbf{oriented PCA}
(oPCA): take the top-$R$ eigenvectors of $\bS_c-\gamma\bS_F$, run $k$-means
inside that subspace, and lift the centroids back. It is closed form, one
eigendecomposition per class, and strictly sharper than a global LDA metric:
LDA finds directions separating all classes \emph{on average}; this finds
directions where class $c$ specifically differs from everyone else.

Two alternatives were run as controls. \emph{Discriminative $k$-means} (dkm)
runs ordinary $k$-means in the metric $(I+\gamma\bS_F)^{-1/2}$, which shrinks
directions foreign material occupies. \emph{Conic matching pursuit} (mp)
exploits the fact that adding a unit ray $\bv$ to a cone with residual $r$
changes the score by exactly the column-generation reduced cost
$\Delta\lVert\Pi\rVert^2=(\bv^{\!\top}r)_+^2$, and iterates the one-sided fixed
point $\bv\leftarrow\widehat{\sum_{S}r_\bx-\gamma\sum_{S'}r_\mathbf{a}}$ over
active sets $S=\{\bx:\bv\cdot r_\bx>0\}$. Dropping the $(\cdot)_+$ would turn
this into an eigenproblem, which is wrong for a cone twice over: eigenvectors
are sign-ambiguous and two-sided.

\begin{table}[t]\centering\small
\begin{tabular}{lccc}
\toprule
construction & $\gamma{=}0$ & $\gamma{=}0.25$ & $\gamma{=}0.5$ \\
\midrule
oPCA & 78.98 & 79.38 & \good{79.77} \\
dkm  & 77.92 & 78.48 & 77.97 \\
\midrule
\multicolumn{2}{l}{[$k$-means, $=$ dkm $\gamma{=}0$]} & \multicolumn{2}{r}{77.92} \\
\multicolumn{2}{l}{[whitened NCM]} & \multicolumn{2}{r}{77.67} \\
\multicolumn{2}{l}{[RanPAC bar]} & \multicolumn{2}{r}{80.28} \\
\bottomrule
\end{tabular}
\caption{Discriminative construction, ImageNet-R $T{=}10$ s0, $R{=}32$, raw cone
A-Last (\texttt{exp39}). oPCA decomposes into $+1.06$ for restricting $k$-means
to the class's own principal subspace and a further $+0.79$ for the foreign
term. dkm is flat: shrinking a \emph{metric} is a weaker intervention than
choosing a \emph{subspace}.}
\label{tab:opca}
\end{table}

Table~\ref{tab:opca} is the central positive result. Two caveats are load
bearing. First, the arms are \emph{asymmetric}: later classes avoid earlier
ones, earlier ones never adapt. Within-task negatives make it work from task
$0$, but the effect should---and does---grow with task index, so A-Last is the
honest column and A-Avg flatters it. Second, an independent run of the same
configuration in \texttt{exp41} reads $79.50$ rather than $79.77$ for
oPCA $\gamma{=}0.5$, $R{=}32$; the two files differ in how the foreign pool is
subsampled. We therefore quote \emph{within-file} deltas throughout and treat
$\pm0.3$ as the implementation-level noise floor on a single cell.

\subsection{Axis 4: allocation ($+0.57$)}
\label{sec:alloc}

Every $R$ sweep before this one was global, which conflates two regimes.
ImageNet-R fit rows run $35$ to $308$ per class, an $8.8\times$ imbalance, and
$24.5\%$ of classes have fewer than 64 rows, so at $R{=}64$ a quarter of the
arms silently clamp and the budget is not matched. Setting
\begin{equation}
R_c=\operatorname{clip}(n_c/K,\;R_{\min},\;R_{\max})
\end{equation}
removes an arbitrary constant rather than adding a mechanism.

\begin{table}[t]\centering\small
\begin{tabular}{lccc}
\toprule
allocation & rays/cls & A-Last & A-Avg \\
\midrule
fixed $R{=}32$              & 32.0 & 79.50 & 84.90 \\
$K{=}2$                     & 51.3 & 79.63 & 84.55 \\
$K{=}3$                     & 34.1 & 79.90 & 84.79 \\
$K{=}5$                     & 20.3 & 79.92 & 84.70 \\
$K{=}8$                     & 13.0 & 79.67 & 84.66 \\
$K{=}12$                    & \phantom{0}9.7 & 79.70 & 84.72 \\
$K{=}5$, $R_{\min}{=}24$    & 26.5 & \good{80.07} & 84.88 \\
$K{=}5$, $R_{\min}{=}32$    & 32.9 & 79.83 & 84.98 \\
$K{=}4$, $R_{\min}{=}32$    & 34.6 & 79.78 & 84.82 \\
\midrule
RanPAC bar                  & ---  & 80.28 & 85.13 \\
\bottomrule
\end{tabular}
\caption{Per-class ray budgets, oPCA $\gamma{=}0.5$, $R\in[8,128]$
(\texttt{exp41}). The best cell uses \emph{fewer} rays than fixed $R{=}32$
($0.84\times$ storage) and scores $+0.57$ higher. The curve is flat between
$K{=}3$ and $K{=}12$---a $3.5\times$ range in mean rays for $0.25$ of
accuracy---so the mechanism is the \emph{floor} $R_{\min}$, not the slope.}
\label{tab:alloc}
\end{table}

The most useful reading of Table~\ref{tab:alloc} is the insensitivity. Between
$9.7$ and $34.1$ mean rays per class the answer moves by a quarter of a point.
A conic read-out on these features is not starved for generators; it is starved
for the \emph{right} generators, which is why axis 3 outperforms axis 4 by
$3\times$.

\subsection{Axis 5: the constraint is atom-dependent and tail-driven}
\label{sec:constraint}

Identical atoms, identical metric, identical score functional---only the
constraint on $\bw$ varies. A rank-$R$ subspace costs exactly the same $Rd$
floats, so \texttt{sub} isolates non-negativity and nothing else.

\begin{table}[t]\centering\small
\begin{tabular}{lcccc}
\toprule
atoms & sub & cone & simplex & cone$-$sub \\
\midrule
$k$-means, $R{=}32$      & \good{78.02} & 77.92 & 77.30 & $-0.10$ \\
oPCA $\gamma{=}0.5$, $k5m24$ & 79.17 & \good{79.77} & 79.57 & \good{$+0.60$} \\
\bottomrule
\end{tabular}
\caption{Region primitive at matched storage, ImageNet-R $T{=}10$ s0. Row 1 is
\texttt{exp38}; row 2 is \texttt{exp45}, which reproduces row 1 exactly before
switching atoms. \emph{The sign of the constraint's value depends on the
generator set}, and the published cell used the $k$-means atoms the method
subsequently abandoned (\S\ref{sec:opca}, \S\ref{sec:alloc}).}
\label{tab:constraint}
\end{table}

Row 1 is also weaker than a single number suggests: across the 10 stages of that
same run the cone beats the subspace on $7$ of them (mean $+0.18$, mid-run
stages 2--8 $+0.37$, sign test $p{=}0.17$), and A-Last is the one stage that
favours \texttt{sub}. The cone leads on A-Avg, $83.28$ to $83.10$.

\paragraph{The constraint is always active.} A natural explanation for a null
result---that the unconstrained solution is already non-negative, so the two
primitives are the same object---is false. At oPCA atoms the constraint binds on
$100.0\%$ of query--class pairs, and $43.7\%$ of the unconstrained
least-squares coefficient mass on the \emph{true} class is negative
(\texttt{exp45}). A query from class $c$ is emphatically not a non-negative
combination of class $c$'s own generators.

\paragraph{Why it helps despite hurting on average.} Define
$\mathrm{gap}(\bq,\bV)=\lVert P_{\mathrm{span}\bV}\bq\rVert-\lVert\Pi_{C(\bV)}\bq\rVert\ge0$,
the amount the constraint removes. Averaged over random foreign classes the
differential runs the \emph{wrong} way: $0.0245$ is removed from the true class
against $0.0144$ from a foreign one, $-9.0\%$ of the mean decision margin.

That mean is dominated by classes that were never in the running. Restricting to
the one foreign class that actually competes---the top-scoring non-true class,
selected under the \emph{subspace} ranking so the comparator is chosen
independently of the constraint being tested---the sign reverses: the runner-up
loses $0.0319$ against the true class's $0.0245$, a differential of
$\mathbf{+0.0074}$, or $+6.6\%$ of the margin. Counted as decisions, the
constraint rescues $196$ true-vs-runner-up contests and breaks $91$.
\textbf{Non-negativity is a tail mechanism: it is a net cost on the average class
and a net gain on the contested one.} This is consistent with \texttt{exp40}'s
finding that confusion is extremely concentrated ($29.3\%$ of a class's errors
go to one confuser, median 4 confusers out of 199).

\paragraph{A counterweight that still stands.} On \emph{near-OOD}, the conic
coefficient carries no signal beyond the residual norm. Decompose the score into
$r_\perp=\lVert\bq-\Pi_C\bq\rVert$ and the ``conic-ness'' $\nu$, the fraction of
unconstrained coefficient mass surviving the projection. $r_\perp$ alone gives
AUROC $0.966$ and $\nu$ alone $0.720$ at $R{=}4$, but $\nu$ \emph{conditioned}
on $r_\perp$ gives $0.602$, and a cross-validated classifier on $(r_\perp,\nu)$
beats one on $r_\perp$ alone by $-0.0008$ (\texttt{exp23}, 40 cells; negative at
$R{=}4,8,16$). The constraint earns its place in the closed-set argmax, not in
the open-set score.

\paragraph{The larger effect is still softness, not sign.} Cone versus
\emph{hard nearest-prototype}---the $k{=}1$ member of its own nesting---is an
order of magnitude bigger than cone versus subspace, and grows with $R$:

\begin{table}[t]\centering\small
\begin{tabular}{lccc}
\toprule
$R$ & max-cosine (pm) & conic (NNLS) & $\Delta$ \\
\midrule
4  & 78.85 & 79.50 & $+0.65$ \\
16 & 77.17 & 78.78 & $+1.61$ \\
64 & 73.47 & 78.02 & \good{$+4.55$} \\
\bottomrule
\end{tabular}
\caption{Soft conic mixture versus hard nearest-prototype at identical
$k$-means atoms (\texttt{exp35}). Multi-prototype \emph{degrades} with more
atoms; the cone does not. The conic rule is an \emph{aggregator} that makes a
noisy atom set usable, not a better class descriptor.}
\label{tab:rule}
\end{table}

Tables~\ref{tab:constraint} and~\ref{tab:rule} together give the correct
mechanistic statement, in two parts. \textbf{Most of the conic read-out's value
is soft aggregation over many atoms} ($+4.55$ at $R{=}64$ over hard
nearest-prototype), which a subspace read-out would also capture.
\textbf{Non-negativity adds a smaller, genuinely conic increment ($+0.60$) that
appears only once the atoms are discriminative, and it acts on contested pairs
rather than on the average class.} The two are separable and were conflated in
an earlier draft of this report, which quoted the $k$-means cell as though it
settled the question.

\subsection{Axis 6: calibration fails}

$\lVert\Pi_C\bq\rVert$ grows with a cone's solid angle, so a class with more
spread data scores higher on \emph{every} query. RanPAC's least-squares fit
calibrates this away for free; the cone has no such mechanism. The bias is real
and measurable: the correlation between a class's cone width and its
over-prediction ratio (times predicted / times true) is $+0.395$
(\texttt{exp40}).

Correcting it makes things worse. Subtracting the class's mean score on foreign
stored rays ($s_c\leftarrow s_c-\mu_{\text{bg},c}$, recomputed every stage, so
drift-free) costs $1.57$ at fixed $R{=}32$ and $2.57$ at $K{=}5$
(\texttt{exp41}: $79.50\to77.93$, $79.92\to77.35$). Dividing additionally by a
per-class $\sigma$ estimated from birth-time validation rows is catastrophic
($90.9\to59.8$ on the cell where it was first tried). The bias exists but it is
apparently \emph{informative}: wide cones belong to classes that genuinely
occupy more of the space, and flattening that removes signal along with bias.

Pruning old classes' rays---the natural storage-side companion---is likewise
monotonically harmful: dropping to $0.9$ and $0.8$ of each class's energy costs
$1.39$ and $3.53$ respectively (\texttt{exp42}, $79.92\to78.53\to76.38$).

\section{The recipe}

\begin{algorithm}[t]
\small
\caption{Best raw conic classifier (oPCA-$k5m24$)}
\label{alg:recipe}
\begin{algorithmic}[1]
\State \textbf{state:} scatter $\Sigma\!\leftarrow\!0$, count $n\!\leftarrow\!0$,
       rays $\{\bV_c\}$
\For{task $t=0,\dots,T-1$}
  \For{class $c$ in task $t$}
    \State $\Sigma \mathrel{+}= X_c^{\!\top}X_c$ (row-centred), $n\mathrel{+}=|X_c|$
  \EndFor
  \State $\bar\Sigma \leftarrow \Sigma/n + \delta\tfrac{\operatorname{tr}}{d}I$,
         \ $\delta{=}0.03$
  \State $\bW_t \leftarrow \operatorname{chol}(\bar\Sigma^{-1})$
  \For{class $c$ in task $t$}
    \State $X_w \leftarrow \widehat{X_c\bW_t}$
    \State $F \leftarrow$ other classes' rows in task $t$ \ $\cup$ \ stored rays
           of past classes, subsampled to $2000$; $F_w\leftarrow\widehat{F\bW_t}$
    \State $R_c \leftarrow \operatorname{clip}(|X_c|/5,\,24,\,128)$
    \State $\bU \leftarrow$ top-$R_c$ eigvecs of
           $\bS_{X_w}-0.5\,\bS_{F_w}$ \Comment{oPCA}
    \State $\bV_c \leftarrow \widehat{\,k\text{-means}_{R_c}(X_w\bU)\,\bU^{\!\top}\bW_t^{-1}}$
  \EndFor
  \State \textbf{predict} $\arg\max_c\lVert\Pi_{C(\widehat{\bV_c\bW_t})}\widehat{\bq\bW_t}\rVert$
         \Comment{FISTA NNLS, 500 iters}
\EndFor
\end{algorithmic}
\end{algorithm}

Algorithm~\ref{alg:recipe} is the whole method. It stores $R_c$ unit vectors per
class ($26.5$ on average, $0.84\times$ the fixed-$R{=}32$ budget) plus one
$768\times768$ scatter, revisits no past image, and has three hyperparameters
($\delta$, $\gamma$, $K$) of which only $\gamma$ matters.

\begin{table}[t]\centering\small
\begin{tabular}{lrr}
\toprule
read-out & A-Last & $\Delta$ \\
\midrule
NCM (raw cosine)                       & 71.02 & --- \\
\quad$+$ tied within-class whitening   & 77.67 & $+6.65$ \\
\quad$+$ $k$-means conic hull, $R{=}32$& 77.92 & $+0.25$ \\
\quad$+$ oriented PCA, $\gamma{=}0.5$  & 79.50 & $+1.58$ \\
\quad$+$ adaptive $R_c$ ($K{=}5$, $R_{\min}{=}24$)
                                       & \good{80.07} & $+0.57$ \\
\midrule
RanPAC (bar)                           & 80.28 & $-0.21$ \\
oracle selector (cone $\vee$ RanPAC)   & 83.42 & \\
\bottomrule
\end{tabular}
\caption{The ledger. ImageNet-R $T{=}10$, seed 0. The metric is $73\%$ of the
total gain and is not conic; the conic rule at fixed atoms is $3\%$;
discriminative selection is $17\%$; allocation is $6\%$.}
\label{tab:ledger}
\end{table}

\section{Where the remaining 0.21 lives}

\texttt{exp40} scores only the final stage (the whitener and generators still
accumulate over all $T$ stages, so the state is correct) and runs five
diagnostics.

\paragraph{The deficit is a capacity mismatch, not geometry.} Binning the final
stage by fit rows per class:

\begin{table}[h]\centering\small
\begin{tabular}{lrrrr}
\toprule
fit rows & $n$ test & cone & RanPAC & gap \\
\midrule
0--50    & \phantom{0}239 & 71.97 & 57.74 & \good{$+14.23$} \\
50--100  & 1559 & 78.58 & 75.63 & $+2.95$ \\
100--200 & 3433 & 80.31 & 82.61 & $-2.30$ \\
200+     & \phantom{0}769 & 80.10 & 86.35 & $-6.24$ \\
\bottomrule
\end{tabular}
\end{table}

\noindent Both extremes are $\sim4.5\sigma$. The cone crushes RanPAC where the
ridge column is starved and underfits, and loses where RanPAC's shared Gram
keeps exploiting samples the cone's $R_c$ rays have saturated on. The $0.21$
aggregate is the residue of two large opposite effects.

\paragraph{Recency is not the problem.} Accuracy by birth task shows no trend,
which re-confirms at $R{=}32$ that the birth-metric-plus-re-metrication scheme
is lossless.

\paragraph{Errors are near-misses.} On errors, the median rank of the true class
among all seen classes is $5$; $32.2\%$ of errors are rank $2$ and $55.1\%$ are
within rank $5$. Margin-level fixes are therefore not hopeless in principle.

\paragraph{Confusion is extremely concentrated.} $29.3\%$ of a class's errors go
to its single top confuser and $61.3\%$ to its top three; the median number of
distinct wrong classes per true class is $4$ out of $199$. This is the sharpest
open lead. The oPCA criterion of \S\ref{sec:opca} provably optimises the gap
between own-class and \emph{average} foreign score, but a classification error
happens when \emph{one} foreign class beats the true one. Pooling $\bS_F$
uniformly gives each of the four real confusers $\sim1/200$ of the weight,
diluting the only mechanism that has ever worked by $\sim50\times$. Weighting
$\bS_F$ by a birth-time confusability estimate
$o(c,c')=\mathbb{E}_{\bx\in X_c}\max_{f\in\text{material}(c')}\langle\bx,f\rangle$
is the obvious next experiment; a per-class-uniform arm is needed in between,
because the foreign pool is inhomogeneous (current-task classes contribute
$\sim2160$ raw rows while each past class contributes $\sim26$ rays, a mix that
drifts from $100\%$ rows at task $0$ to $27\%$ at task $9$). This is implemented
in \texttt{exp43} and \textbf{not yet run}.

\section{The fusion ceiling, and why it is not the method}
\label{sec:fusion}

The $2\times2$ of (cone right/wrong) $\times$ (RanPAC right/wrong) at the final
stage is $\{$both $4607$, cone-only $188$, RanPAC-only $210$, neither $995\}$,
so a per-query oracle selector reaches $83.42$ against RanPAC's $80.28$: the
two channels are genuinely complementary, and the complementarity \emph{grows}
over stages while their score correlation \emph{falls} from $0.81$ to $0.50$.

That headroom is not extractable. Three structurally different combiners---a
global $\beta$, a per-sample gate, and an EMA-anchored $\beta$---land within
$0.25$ of one another and capture $8$--$26\%$ of the oracle, reaching $81.03$
A-Last / $85.99$ A-Avg at $R{=}64$ (\texttt{exp35}). Notably the conic rule's
\emph{fused} value also scales with rays ($+0.08/+0.45/+0.75$ over fused
multi-prototype at $R{=}4/16/64$), mirroring Table~\ref{tab:rule}. We report
these numbers for completeness and exclude them from the method: a fused
headline moves by $0.77$ under a $0.09$ feature perturbation through $\beta$
selection alone, which makes it useless as an attribution instrument.

\section{Negative results}

\paragraph{Lifting the cone.} RanPAC's random ReLU projection is worth $+9.26$
over NCM---larger than every other component combined---and every cone above
operates \emph{beneath} it. Moving generators and NNLS into
$\mathrm{relu}(\phi P)$ gains nothing ($-0.07$ for the cone, $+0.20$ for
max-cosine on ImageNet-R; $53.93\to53.87$ on ImageNet-A). The reason is
measurable: the lift \emph{doubles} mean pairwise cosine, $0.207\to0.429$, and
does not improve with $M$ ($0.4285$ at $M{=}2000$, $0.4289$ at $M{=}10000$).
Crowding every point into the non-negative orthant is harmless to a ridge, which
can build discriminative directions anywhere, and harmful to any reader that
sees only angles. The lift belongs to the ridge, not to the representation.

\paragraph{The $H$-representation.} Replacing the generator description with a
facet description fixes the monotonicity of coverage in $R$ but fails three
pre-registered gates. The decisive one is aggregation: min-over-facets
($L^\infty$) $<$ softmin $<$ sum-of-squares ($L^2$), monotone in \emph{how
little the rule resembles a conjunction}. Conic membership \emph{is} a
conjunction, so the conic reading is structurally locked into the worst pooling.

\paragraph{Multi-layer features.} Concatenating or pooling intermediate ViT
blocks strictly hurts every arm (\texttt{exp37}: \texttt{cls} $80.37$,
\texttt{cls+gap} $80.08$, multi-layer \texttt{cls} $79.77$, multi-layer
\texttt{gap} $77.20$). The final CLS token is the right read-out point.

\paragraph{Test-time adaptation.} Transductive refinement of the cone using the
test batch moves nothing: mean-shift $79.80\to79.77$, residual-shift
$79.80\to79.80$, and $\pm0.02$ for multi-prototype. Ray pruning during TTA costs
$1.35$ for the cone and $1.15$ for multi-prototype.

\paragraph{Between-class subspace overlap is not the failure mode.}
\texttt{exp29} correlates per-class error against candidate geometric
quantities. Between-class subspace affinity does not predict it
($\rho={+}0.037$ at rank 32; every trivial control, including sample count at
$\rho={-}0.347$, predicts better). \emph{Within}-class coverage does: the mean
own-subspace energy $\lVert P_c\bq\rVert$ of a class's held-out rows has
$\rho={-}0.585$ and strengthens with rank. \textbf{Classes fail because their
own test rows fall outside their own training span---not because classes
collide.} This is the diagnostic that motivates \texttt{exp44}, an effective-rank
shaping loss on task 0 ($L=\mathrm{CE}+\lambda_w\,\overline{\erank(\bS_c)}
-\lambda_p\,\erank(\bS_{\text{pool}})$, with the pooled term required because
the head whitens by $\bS_w$ and would otherwise destroy any uniformly low-rank
structure). Feature caches exist; \textbf{results do not yet}.

\section{The regime where the raw cone already wins}

Everything above is ImageNet-R, which has $108$ fit rows per class on average.
The class-size table predicts that the cone should win outright on datasets that
sit in its left-hand bins, and it does. On ImageNet-A ($\sim18$ rows/class) at
$R{=}4$ with a task-0 whitener, the raw cone scores $57.13$ A-Last against
RanPAC's $55.27$ (\texttt{exp33}), a $+1.86$ win with no fusion and no
calibration; the max-cosine arm on the same atoms scores $52.53$, so the conic
rule contributes $+4.60$ of it. CUB-200 ($\sim27$ rows/class) is the other
predicted cell.

This is the honest positioning of the primitive. A per-class generative region
model is not a better classifier than a jointly fitted discriminative head when
data is plentiful---it is naive Bayes against LDA, and it loses for the usual
reason. It is a better classifier when the discriminative head cannot estimate
its own parameters, which in class-incremental learning is the common case at
the tail of the class distribution and the whole case on few-shot benchmarks.

\section{Limitations}
\label{sec:limits}

\paragraph{Seeds.} \emph{Every read-out number in this report is seed 0 only.}
RanPAC alone spans $80.28/80.55/80.38$ A-Last across three seeds on ImageNet-R
$T{=}10$---a $0.27$ range, the same order as the allocation effect
(\S\ref{sec:alloc} $=+0.57$) and larger than the final deficit ($0.21$). The
oPCA effect ($+1.58$) is the only single-axis result comfortably outside seed
noise. No other claim here should be treated as established until it is seeded.

\paragraph{One cell.} Almost every ablation is ImageNet-R $T{=}10$. The
allocation, calibration, pruning and constraint results have not been rerun at
$T{=}20/50$ or on CIFAR-100, CUB-200, ImageNet-A. The $\pm0.3$ cross-file
discrepancy documented in \S\ref{sec:opca} bounds how much a single cell can be
trusted.

\paragraph{Contaminated reference columns.} An earlier NCM implementation read
class means off the first arm's generators, which made the NCM column vary with
$R$ ($76.72/78.02/78.57/78.85$ at $R{=}4/16/32/64$---a quantity that cannot
depend on $R$). \texttt{exp39} onward computes a true class mean. Any NCM
number quoted from \texttt{exp35} or earlier is unreliable; the $77.67$ used in
Table~\ref{tab:ledger} is the corrected one.

\paragraph{Scope.} One backbone, one adaptation recipe. The geometric claims
concern pre-trained ViT features in this regime and should not be read as claims
about conic parameterisations generally.

\section{Conclusion}

The best raw conic classifier we can build reaches $80.07$ A-Last on
ImageNet-R $T{=}10$ against a RanPAC bar of $80.28$, from $71.02$ for
nearest-class-mean, storing $26.5$ vectors per class and revisiting no image.
The recipe is short: whiten by the accumulated tied within-class scatter, choose
generators by oriented PCA against the legally available foreign material
($\gamma{=}0.5$), allocate $R_c\propto n_c$ with a floor of 24, and score with a
properly converged NNLS projection.

The attribution is less flattering than the headline. Of the $9.05$ points, the
metric is $6.65$ and is not conic; discriminative \emph{selection} is $1.58$;
the conic \emph{rule} at fixed $k$-means atoms is $0.25$, of which
non-negativity contributes nothing at those atoms and $+0.60$ at the
discriminative ones. What survives as a genuine property of
the primitive is narrower and, we think, more useful: a soft mixture read-out
degrades gracefully in the number of atoms where hard nearest-prototype
collapses ($+4.55$ at $R{=}64$), it produces an absolute membership score, and
it beats a discriminative head precisely where that head runs out of samples
($+14.23$ below 50 rows per class, $+1.86$ overall on ImageNet-A). The remaining
$0.21$ on ImageNet-R is not a geometry problem; it is a capacity mismatch on
large classes, and the concentration of confusions---four confusers out of
199---says the next thing to try is a confusability-weighted $\bS_F$, not
another region primitive.

\end{document}
